% ---------------------------------------------------------------
% Preamble — biên dịch bằng XeLaTeX (Overleaf: Menu > Compiler > XeLaTeX)
% ---------------------------------------------------------------
\usepackage[a4paper,margin=2.5cm]{geometry}
\usepackage{fontspec}
\usepackage{polyglossia}
\setdefaultlanguage{vietnamese}
% TeX Gyre Termes = font tương đương Times New Roman. Nạp theo tên file để chạy được cả trên tectonic lẫn Overleaf/TeX Live
\setmainfont{texgyretermes}[Extension=.otf, UprightFont=*-regular, BoldFont=*-bold,
  ItalicFont=*-italic, BoldItalicFont=*-bolditalic]
\setmonofont{texgyrecursor}[Extension=.otf, UprightFont=*-regular, BoldFont=*-bold,
  ItalicFont=*-italic, BoldItalicFont=*-bolditalic, Scale=0.95]

\usepackage{setspace}
\onehalfspacing
\setlength{\emergencystretch}{3em}  % giảm dòng tràn lề khi có tên file dài
\usepackage{microtype}
\usepackage{amsmath,amssymb}
\usepackage{graphicx}
\usepackage{adjustbox}
\usepackage{booktabs,tabularx,longtable,multirow}
\usepackage{array}
\usepackage[table]{xcolor}
\usepackage{caption}
\captionsetup{font=small,labelfont=bf}
\usepackage{float}
\usepackage{enumitem}
\setlist{nosep,leftmargin=*}
\usepackage{tikz}
\usetikzlibrary{arrows.meta,positioning,shapes,fit,backgrounds,calc}
\usepackage{tcolorbox}
\tcbuselibrary{skins,breakable}
\usepackage{listings}
\usepackage{titlesec}
% Tiêu đề chương gọn (mặc định của lớp report để trống nửa trang)
\titleformat{\chapter}[hang]{\LARGE\bfseries}{\thechapter.}{0.6em}{}
\titlespacing*{\chapter}{0pt}{-10pt}{14pt}
\usepackage{fancyhdr}
\usepackage[hidelinks]{hyperref}

% --- Header/footer ---
\pagestyle{fancy}
\fancyhf{}
\fancyhead[L]{\small CSE703093 — An toàn phần mềm}
\fancyhead[R]{\small Bài tập lớn}
\fancyfoot[C]{\thepage}

% --- Màu ---
\definecolor{bad}{HTML}{C62828}
\definecolor{good}{HTML}{2E7D32}
\definecolor{note}{HTML}{1565C0}
\definecolor{codebg}{HTML}{F6F8FA}

% --- Macro ---
% \TODO{...}: chỗ cần điền kết quả thật từ results/ — xóa hết trước khi nộp
\newcommand{\TODO}[1]{\textcolor{bad}{\textbf{[TODO: #1]}}}
\newcommand{\cwe}[2]{\textbf{CWE-#1} (#2)}
% \code{...}: cho phép xuống dòng sau "_" trong tên hàm, tên file dài
\newcommand{\codeus}{\textunderscore\allowbreak}
\DeclareRobustCommand{\code}[1]{\texttt{\let\_\codeus #1}}
\newcommand{\req}[1]{\textsf{#1}}

% Hộp ghi chú
\newtcolorbox{ghichu}[1][]{colback=note!5,colframe=note!60,boxrule=0.5pt,
  arc=2pt,left=6pt,right=6pt,top=4pt,bottom=4pt,fonttitle=\bfseries\small,title={#1}}

% --- Listings (giữ mã nguồn ASCII để tránh lỗi mã hóa) ---
\lstdefinestyle{base}{
  basicstyle=\ttfamily\footnotesize,
  backgroundcolor=\color{codebg},
  frame=single, rulecolor=\color{black!15},
  numbers=left, numberstyle=\tiny\color{black!50}, numbersep=6pt,
  keywordstyle=\color{note}\bfseries,
  commentstyle=\color{black!55}\itshape,
  stringstyle=\color{good!70!black},
  showstringspaces=false, breaklines=true, columns=fullflexible,
  xleftmargin=1.5em, framexleftmargin=1.5em, captionpos=b
}
\lstset{style=base}
\lstdefinelanguage{log}{morecomment=[l]{\#}}
\lstdefinestyle{logstyle}{style=base,language=log,numbers=none,backgroundcolor=\color{black!3}}
\renewcommand{\lstlistingname}{Đoạn mã}

% --- TikZ styles ---
\tikzset{
  st/.style={draw,circle,minimum size=13mm,align=center,font=\scriptsize,fill=white},
  init/.style={st,fill=note!15},
  viol/.style={st,fill=bad!25,draw=bad},
  box/.style={draw,rounded corners=2pt,minimum height=9mm,align=center,font=\small,fill=white},
  arr/.style={-{Stealth[length=2mm]},semithick},
}
